\documentclass[10pt]{article}
\providecommand{\CRevisionRound}{2}
\usepackage[margin=0.76in]{geometry}
\usepackage{fontspec}
\setmainfont{TeX Gyre Pagella}
\setsansfont{TeX Gyre Heros}
\setmonofont{Latin Modern Mono}
\usepackage{microtype,amsmath,amssymb,amsthm,mathtools,booktabs,array,enumitem}
\usepackage[hidelinks]{hyperref}
\pdfvariable suppressoptionalinfo 767
\pdfvariable trailerid {[<C2842026090200000000000000000000><C2842026090200000000000000000000>]}
\setlength{\parindent}{0pt}
\setlength{\parskip}{4pt}
\setlength{\emergencystretch}{3em}
\setlist{nosep,leftmargin=*}
\newtheorem{theorem}{Theorem}
\newtheorem{lemma}{Lemma}
\newtheorem{proposition}{Proposition}
\newcommand{\R}{\mathbb R}
\newcommand{\J}{\mathsf J}
\newcommand{\Id}{\mathsf I}
\newcommand{\diag}{\operatorname{diag}}
\newcommand{\Span}{\operatorname{span}}
\title{The Thomson Polygon from Its Cartesian Hessian:\\
Exact Fourier Blocks and the Sharp Linear Stability Threshold}
\author{HCS Research Program}
\date{2 September 2026}

\begin{document}
\maketitle

\begin{abstract}
For identical positive point vortices on a planar regular $N$-gon, we derive
the rotating relative equilibrium directly from the logarithmic Hamiltonian
and reduce its raw Cartesian Hessian by a vertexwise radial--tangential
discrete Fourier transform.  If
$c=\Gamma/(4\pi R^2)$ and $q_m=m(N-m)$, the $m$th Hamiltonian block is
\[
 c\begin{pmatrix}0&q_m\\-[2(N-1)-q_m]&0\end{pmatrix}.
\]
Root-of-unity orthogonality proves this formula for every $N\ge3$.
Consequently the symmetry-reduced polygon is linearly elliptic for
$3\le N\le6$, linearly degenerate only in the conjugate modes $m=3,4$ when
$N=7$, and hyperbolic for every $N\ge8$.
\ifnum\CRevisionRound>0
We separate rotation, scale, translations, and the centered first-harmonic
complement, and close the $\Gamma$, $R$, collision, and low-$N$ faces.
\fi
\ifnum\CRevisionRound>1
An independent checker reconstructs every raw $2N\times2N$ Hessian for
$3\le N\le64$; exact symbolic and mutation audits close the release.
\fi
The heptagon statement is strictly linear.  The result is a self-contained
reconstruction of a classical owner theorem, not a claim of literature
priority or nonlinear heptagon stability.
\end{abstract}

\section{Frozen Hamiltonian and relative equilibrium}

Let $z_j=(x_j,y_j)\in\R^2$ and give every vortex the same circulation
$\Gamma>0$.  We freeze the convention
\begin{equation}
 H(z)=-\frac{\Gamma^2}{2\pi}\sum_{j<k}\log|z_j-z_k|,
 \qquad
 \Gamma\dot z_j=\J\nabla_jH,
 \qquad
 \J=\begin{pmatrix}0&1\\-1&0\end{pmatrix}.                 \label{eq:model}
\end{equation}
Thus $-\J$ generates counterclockwise rotation.  Take
\begin{equation}
 a_j=R(\cos\theta_j,\sin\theta_j),\qquad
 \theta_j=\frac{2\pi j}{N},\qquad N\ge3,\quad R>0.          \label{eq:polygon}
\end{equation}
Thomson~\cite{Thomson1883} and Havelock~\cite{Havelock1931} are cited as
classical owners of the polygon and its linear-stability question.
Cabral and Schmidt~\cite{CabralSchmidt2000} provide later stability context,
and Celli, Lacomba, and P\'erez-Chavela~\cite{Celli2011} document polygonal
relative-equilibrium lineage.  No calculation or proof below is outsourced
to these sources, and no literature-priority claim is made.

\begin{lemma}[Angular velocity]
The configuration \eqref{eq:polygon} is a relative equilibrium with
\begin{equation}
 \Omega=\frac{\Gamma(N-1)}{4\pi R^2}.                       \label{eq:omega}
\end{equation}
For
$G=H+(\Gamma\Omega/2)\sum_j|z_j|^2$, one has $\nabla G(a)=0$.
\end{lemma}

\begin{proof}
Differentiating \eqref{eq:model} gives
\[
 \nabla_jH=-\frac{\Gamma^2}{2\pi}
 \sum_{k\ne j}\frac{z_j-z_k}{|z_j-z_k|^2}.
\]
At $j=0$, each summand inside the sum has radial component $1/(2R)$;
tangential components cancel between $k$ and $N-k$.  Rotational covariance
therefore gives
\[
 \sum_{k\ne j}\frac{a_j-a_k}{|a_j-a_k|^2}
 =\frac{N-1}{2R^2}a_j.
\]
Hence $\nabla_jH(a)=-\Gamma\Omega a_j$ with \eqref{eq:omega}, proving both
claims and the sign of rotation under \eqref{eq:model}.
\end{proof}

The rotating-frame linearization is
\begin{equation}
 \mathcal L=\J_N\,\Gamma^{-1}D^2G(a),
 \qquad \J_N=\diag(\J,\ldots,\J).                           \label{eq:linearization}
\end{equation}
All stability statements below refer to \eqref{eq:linearization}; they do not
silently change the Hamiltonian or clock.

\section{Raw Cartesian Hessian and exact Fourier blocks}

For $d\ne0$, put
\[
 A(d)=\frac{\Id}{|d|^2}-\frac{2dd^T}{|d|^4}.
\]
The pairwise Cartesian blocks are
\begin{equation}
 D^2_{jj}H=-\frac{\Gamma^2}{2\pi}A(d_{jk}),
 \qquad
 D^2_{jk}H=+\frac{\Gamma^2}{2\pi}A(d_{jk}),                \label{eq:rawhessian}
\end{equation}
and the augmentation adds $\Gamma\Omega\Id$ to each diagonal block.
Equation \eqref{eq:rawhessian}, rather than the final closed formula, is the
starting point of the independent checker.

Let $Q_j$ have the radial and counterclockwise tangential unit vectors at
$a_j$ as columns.  Set
\begin{equation}
 c=\frac{\Gamma}{4\pi R^2},\qquad q_m=m(N-m),\qquad
 \sigma_m=2(N-1)-q_m.                                      \label{eq:notation}
\end{equation}

\begin{theorem}[Complete DFT block]
In the local radial--tangential DFT, for every $m=0,\ldots,N-1$,
\begin{align}
 \widehat{\Gamma^{-1}D^2G}_m
   &=c\diag(\sigma_m,q_m),                                  \label{eq:hblock}\\
 \mathcal L_m
   &=c\begin{pmatrix}0&q_m\\-\sigma_m&0\end{pmatrix},
 &\mathcal L_m^2&=-c^2q_m\sigma_m\Id.                      \label{eq:lblock}
\end{align}
\end{theorem}

\begin{proof}
For $k\ne0$, let $\theta=\theta_k$ and $d=a_0-a_k$.  Since
$|d|^2=4R^2\sin^2(\theta/2)$, substitution in $A(d)$ gives
\[
 A(d)=\frac1{4R^2\sin^2(\theta/2)}
 \begin{pmatrix}\cos\theta&\sin\theta\\
                 \sin\theta&-\cos\theta\end{pmatrix}.
\]
Multiplication by $Q_k$ turns the displayed matrix into $\diag(1,-1)$.
Thus the local off-diagonal block in the first block row is
\[
 B_k=\frac{c}{1-\cos\theta_k}\diag(1,-1).
\]
If $A_N=\sum_{k=1}^{N-1}(1-\cos\theta_k)^{-1}$, the raw diagonal block is
$B_0=c\diag(2(N-1)-A_N,A_N)$.  Pairing $k$ with $N-k$ cancels the sine
parts of the DFT and yields
\begin{equation}
 \widehat{\Gamma^{-1}D^2G}_m
 =c\diag(2(N-1)-S_m,S_m),
 \quad
 S_m=\sum_{k=1}^{N-1}
 \frac{1-\cos(m\theta_k)}{1-\cos\theta_k}.                 \label{eq:rootsum}
\end{equation}

For $\zeta=e^{2\pi i/N}$,
\[
 \frac{1-\cos(m\theta)}{1-\cos\theta}
 =\left|\sum_{r=0}^{m-1}e^{ir\theta}\right|^2.
\]
Root-of-unity orthogonality therefore gives
\[
 \sum_{k=0}^{N-1}\left|\sum_{r=0}^{m-1}\zeta^{kr}\right|^2=Nm.
\]
The $k=0$ term is $m^2$, so $S_m=Nm-m^2=q_m$.  This proves
\eqref{eq:hblock}.  The local rotations preserve $\J$; multiplying by $\J$
and squaring proves \eqref{eq:lblock}.
\end{proof}

\section{The sharp linear threshold}

\begin{theorem}[Thomson polygon trichotomy]
After center, angular-impulse, and rotation reduction:
\begin{enumerate}
\item every block is semisimple elliptic for $3\le N\le6$;
\item for $N=7$, exactly $m=3,4$ are nonzero nilpotent blocks;
\item every $N\ge8$ has a real hyperbolic pair.
\end{enumerate}
The middle statement is linear degeneracy only; it does not assert nonlinear
stability or instability of the heptagon.
\end{theorem}

\begin{proof}
For $m\ne0$, $q_m>0$, so \eqref{eq:lblock} is elliptic, nilpotent, or
hyperbolic according as $\sigma_m$ is positive, zero, or negative.  The
maximum of $q_m$ is $\lfloor N^2/4\rfloor$.  Hence the least sign is
\[
 2(N-1)-\lfloor N^2/4\rfloor.
\]
It is positive for $N=3,4,5,6$ and zero at $N=7$, where
$m(7-m)=12$ only for $m=3,4$.  For $N=2k\ge8$,
$k^2-(4k-2)>0$; for $N=2k+1\ge9$, $k(k+1)-4k=k(k-3)>0$.
Thus the maximal-$q_m$ block is hyperbolic from eight onward.  The symmetry
directions removed in the statement are identified explicitly below in the
revised manuscript and do not alter these signs.
\ifnum\CRevisionRound>1
The labels $m=3,4$ at $N=7$ are conjugate.  On their real four-dimensional
isotypic component, the zero eigenvalue has algebraic multiplicity four and
geometric multiplicity two.  This refines the linear degeneracy statement but
does not supply a nonlinear conclusion.
\fi
\end{proof}

\begin{center}
\small
\begin{tabular}{@{}llll@{}}
\toprule
$N$ & controlling $\sigma_m$ & reduced linear type & strict boundary\\
\midrule
$3$--$6$ & positive & elliptic & semisimple\\
$7$ & zero at $m=3,4$ & degenerate & no nonlinear claim\\
$8$ and above & negative for some $m$ & hyperbolic & real pair\\
\bottomrule
\end{tabular}
\end{center}

\ifnum\CRevisionRound>0
\section{Symmetry slice and parameter faces}

The full Fourier ledger prevents symmetry modes from being mislabeled as
instabilities.  At $m=0$,
\begin{equation}
 \mathcal L_0=c\begin{pmatrix}0&0\\-2(N-1)&0\end{pmatrix}. \label{eq:mzero}
\end{equation}
Uniform tangential displacement is rotation; uniform radial displacement is
its scale generalized vector.  Fixing angular impulse and quotienting
rotation remove \eqref{eq:mzero}.

The conjugate first harmonics contain the translation plane.  If local
coordinates are $\xi_j,\eta_j$ and $w_j=\xi_j+i\eta_j$, then the center
variation is
\[
 \delta\!\left(\sum_jz_j\right)=\sum_je^{i\theta_j}w_j.
\]
Fixing the center removes this complex translation plane.  It does not remove
the entire first-harmonic isotypic component: because
$q_1=\sigma_1=N-1$, the Hessian there is the scalar $c(N-1)\Id$ and its
centered complementary plane remains elliptic with frequency
$c(N-1)=\Omega$.  The centered, fixed-impulse rotational slice therefore
removes exactly the intended Euclidean and scale directions while retaining
every genuine shape degree of freedom.

The physical parameter faces are also explicit.
\begin{itemize}
\item $R=0$ is a logarithmic collision and is excluded.  As $R\to\infty$,
all frequencies vanish like $R^{-2}$ while the sign atlas remains fixed.
\item $\Gamma=0$ degenerates the weighted symplectic form; it is recorded only
as a zero-velocity limit.  Common $\Gamma<0$ reverses time and preserves all
stability signs.
\item $N=1$ is trivial.  After centering, $N=2$ has no reduced shape degree of
freedom, so the theorem begins at $N=3$.
\item $N=7,m=3,4$ is an exact nonzero nilpotent face.  A higher-order
nonlinear conclusion would require a different theorem and is not inferred.
\end{itemize}

\section{Independent executable reconstruction}

The finite certificate tests the derivation without pretending to prove it.
The producer emits \eqref{eq:notation}--\eqref{eq:lblock} as exact integer
cells.  The checker shares no producer function: for each $N=3,\ldots,64$ it
constructs \eqref{eq:rawhessian} in global Cartesian coordinates, verifies
the augmented equilibrium, checks the complete $2N\times2N$ symmetry, rotates
every block by the $Q_j$, verifies cyclicity, performs the DFT, and applies the
raw Hessian and Hamiltonian linearization to explicit rotation, scale,
translation, and centered first-harmonic vectors.

SymPy separately differentiates and transforms exact raw Hessians for
$N=3,4,6$ and verifies the same slice actions.  Its root-sum audit covers every
one of the 2,077 archived mode cells by coefficient counting rather than by
inserting the closed value.  Two unrelated fresh output paths must reproduce
the canonical JSON byte-for-byte.  The checker rejects duplicate keys,
nonstandard constants, unknown or missing fields, bool/type confusion,
reordered or semantically duplicated rows, and all altered contract values,
even when the payload hash is repaired.
\fi

\ifnum\CRevisionRound>1
\section{Receipt, ownership, and Route-A boundary}

The deterministic receipt contains 2,077 exact mode rows, 62 polygon rows, 64
exact $\Gamma/R^2$ scale cells, seven symmetry-slice rows, and eight
boundaries.  The independent raw checker passes 65,655 assertions, the
symbolic reconstruction passes 4,585 identities, the two-path replay is
byte-exact, and 76/76 hostile specimens are rejected.  The latter comprise
73 repaired-hash attacks plus stale-hash, raw duplicate-key, and nonstandard
constant controls.  The evidence SHA-256 begins
\texttt{4fed9820df14c399} and ends \texttt{8123fcdfb39db9462a53}.

These counts are implementation evidence.  The arbitrary-$N$ theorem rests
on \eqref{eq:rootsum} and the exact maximum of $m(N-m)$, not on the cutoff at
$N=64$.  The source audit likewise separates lineage from proof: classical
owners receive explicit attribution, while the workspace increment is only
the convention-frozen synthesis and independently executable closure.

The honest Route-A evaluation is
\[
 (\mathrm{A0\_FAIL},\mathrm{A1\_WEAK},\mathrm{A2\_FAIL},
   \mathrm{A3\_FAIL},\mathrm{A4\_FAIL}),
 \qquad \mathrm{ROUTE\_A\_REJECTED}.                        \label{eq:route}
\]
The polygons provide a natural relative-periodic family, justifying only the
weak A1 label.  Their clock varies continuously with $\Gamma/R^2$; one
symmetric family is not an isolated primitive-orbit census.  There is no
rational-prime carrier, logarithmic-prime clock, target determinant, target
divisor, functional equation, target zero match, Hilbert--P\'olya operator, or
same-clock quantum lift.  Route B is not authorized.  The scope literal is
\texttt{NO\_BAD\_EULER\_OR\_ROOT\_NUMBER}.

The strict nonclaims are therefore substantive: no nonlinear stability result
for the Thomson heptagon, no proof by finite enumeration, no literature
priority claim, and no arithmetic or spectral-target bridge.
\fi

\section{Conclusion}

The logarithmic pair Hessian and one root-of-unity sum determine every
Fourier block of the equal-vortex polygon.  Once rotation, scale, and
translations are handled explicitly, the sign
$2(N-1)-m(N-m)$ yields the exact reduced linear transition at seven and eight
vortices.
\ifnum\CRevisionRound>0
The parameter ledger prevents zero circulation, collision, infinite radius,
or the first harmonics from being folded into the shape theorem.
\fi
\ifnum\CRevisionRound>1
The executable audit reconstructs, rather than assumes, the Cartesian-to-DFT
step.  Its strongest limitation is intentional: linear degeneracy at $N=7$
is not promoted to a nonlinear statement, and classical reconstruction is not
presented as new literature.
\fi

\begin{thebibliography}{9}
\bibitem{Thomson1883}
J. J. Thomson, \emph{A Treatise on the Motion of Vortex Rings}, Macmillan,
London (1883).

\bibitem{Havelock1931}
T. H. Havelock, ``The stability of motion of rectilinear vortices in ring
formation,'' \emph{Philosophical Magazine} \textbf{11} (70), 617--633 (1931),
\href{https://doi.org/10.1080/14786443109461714}{doi:10.1080/14786443109461714}.

\bibitem{CabralSchmidt2000}
H. E. Cabral and D. S. Schmidt, ``Stability of relative equilibria in the
problem of $N+1$ vortices,'' \emph{SIAM J. Math. Anal.} \textbf{31} (2),
231--250 (2000),
\href{https://doi.org/10.1137/S0036141098302124}{doi:10.1137/S0036141098302124}.

\bibitem{Celli2011}
M. Celli, E. A. Lacomba, and E. P\'erez-Chavela, ``On polygonal relative
equilibria in the $N$-vortex problem,'' \emph{J. Math. Phys.} \textbf{52},
103101 (2011),
\href{https://doi.org/10.1063/1.3646115}{doi:10.1063/1.3646115}.
\end{thebibliography}

\end{document}
